% ==========================================================================
%  Chapter 79 -- 16-story L-shaped RC building multiscale seismic analysis
% ==========================================================================

\chapter{Ejemplo avanzado: Edificio en L de 16 niveles con análisis
         multiescala FE\textsuperscript{2}}
\label{ch:multiscale-16storey}

\paragraph{Normative status.}
Este capítulo se conserva como registro histórico del ejemplo de aplicación.
La nomenclatura y la semántica normativas del módulo multiescala vigente se
definen en el Chapter~\ref{ch:multiscale-publication-audit}.

% ------------------------------------------------------------------
\section{Introducción y motivación}
\label{sec:ms16-intro}

Los edificios altos de concreto reforzado (CR) con irregularidades en
altura presentan concentraciones de daño en las zonas de transición de
sección, fenómeno que los modelos de fibra convencionales capturan solo
parcialmente.  En este capítulo se demuestra la capacidad completa del
módulo multiescala FE\textsuperscript{2} de \texttt{fall\_n} para analizar
un edificio en L de 16 niveles con tres zonas de sección diferente,
sometido al sismo de Tohoku 2011 (estación MYG004) en las tres
componentes (NS, EW, UD).

El objetivo es triple:
\begin{enumerate}
  \item \textbf{Validar la transición automática} del análisis global
        con fibras a submodelos 3D No~Lineales (Ko-Bathe Concrete 3D)
        cuando se alcanza la fluencia del acero.
  \item \textbf{Demostrar el acoplamiento FE\textsuperscript{2} doble}
        (two-way staggered) con convergencia de Frobenius y relajación
        constante, aplicado a múltiples submodelos con secciones
        distintas.
  \item \textbf{Capturar la evolución de fisuras} en los submodelos y
        cuantificar su efecto en la respuesta global (desplazamientos
        de cubierta, curvas de histéresis, envolvente de derivas).
\end{enumerate}

% ------------------------------------------------------------------
\section{Hipótesis del modelo}
\label{sec:ms16-hypotheses}

Se plantean las siguientes hipótesis de trabajo:
\begin{description}
  \item[H1 --- Localización del daño:]
    Debido a la reducción brusca de sección entre pisos 5--6 y 11--12,
    el daño se concentrará en las columnas de borde de la re-entrante
    del plano en L, entre los pisos 5 y 6 (transición de $0{.}50\times
    0{.}50$\,m a $0{.}40\times 0{.}40$\,m).
  \item[H2 --- Acumulación progresiva de fisuras:]
    Los submodelos 3D registrarán una apertura de fisura creciente
    con el número de ciclos de carga, alcanzando valores máximos
    durante los pulsos de mayor PGA ($t\approx 60$--$90$\,s).
  \item[H3 --- Efectividad del acoplamiento FE\textsuperscript{2}:]
    El esquema staggered de dos vías convergerá en menos de 6
    iteraciones por paso, con norma de Frobenius relativa menor al
    3\,\%.
  \item[H4 --- Compatibilidad de irregularidad vertical:]
    Las columnas del rango superior ($0{.}30\times 0{.}30$\,m,
    $f'_c=21$\,MPa) experimentarán mayor deformación unitaria pero
    menor fuerza axial, resultando en un índice de daño menor al de
    las columnas de transición.
\end{description}

% ------------------------------------------------------------------
\section{Descripción del modelo}
\label{sec:ms16-model}

\subsection{Geometría}

El edificio tiene planta en forma de L con 4 vanos en $x$ (5\,m cada
uno $=$ 20\,m totales) y 3 vanos en $y$ (4\,m cada uno $=$ 12\,m
totales), con una esquina recortada ($2\times 1$ vanos en la zona
superior derecha).  Se consideran 16 niveles con entrepiso uniforme
$h_e = 3{.}2$\,m, para una altura total de $H = 51{.}2$\,m.

\subsection{Irregularidad en altura --- secciones por rango}

\begin{table}[htbp]
  \centering
  \caption{Secciones de columna por rango de pisos.}
  \label{tab:ms16-col-sections}
  \begin{tabular}{lccccc}
    \toprule
    Rango  & Pisos & $b\times h$ [m] & $f'_c$ [MPa]
           & $E_c$ [MPa] & Grupo \\
    \midrule
    Inferior & 1--5   & $0.50\times0.50$ & 35 & 27\,806 & \texttt{ColumnsLower} \\
    Medio    & 6--11  & $0.40\times0.40$ & 28 & 24\,870 & \texttt{ColumnsMid}   \\
    Superior & 12--16 & $0.30\times0.30$ & 21 & 21\,538 & \texttt{ColumnsUpper} \\
    \bottomrule
  \end{tabular}
\end{table}

Las vigas son uniformes en todos los niveles: $0{.}30\times 0{.}60$\,m,
$f'_c = 25$\,MPa.  Todas las barras longitudinales tienen $f_y = 420$\,
MPa y relación de endurecimiento $b = 0{.}01$.

\subsection{Acción sísmica}

Se emplea el registro del sismo de Tohoku 2011-03-11 ($M_w$ 9.0),
estación K-NET MYG004 (Tsukidate, Miyagi), cerca del campo cercano.
Se aplican las tres componentes simultáneamente:
\begin{itemize}
  \item \textbf{NS} $\rightarrow$ dirección $x$
  \item \textbf{EW} $\rightarrow$ dirección $y$
  \item \textbf{UD} $\rightarrow$ dirección $z$
\end{itemize}
con escala $\alpha = 1{.}0$ y ventana de análisis de 275\,s (descartando
los primeros 25\,s de ruido previo al evento).

\subsection{Parámetros de análisis dinámico}

\begin{table}[htbp]
  \centering
  \caption{Parámetros del análisis dinámico.}
  \label{tab:ms16-dyn-params}
  \begin{tabular}{lcc}
    \toprule
    Parámetro & Valor & Unidad \\
    \midrule
    Paso de tiempo $\Delta t$        & 0.02  & s \\
    Período fundamental $T_1$        & 1.60  & s \\
    Tercer modo $T_3$                & 0.40  & s \\
    Factor de amortiguamiento $\xi$  & 5     & \% \\
    Densidad $\rho$                  & $2.4\times10^{-3}$ & MN\,s$^2$/m$^4$ \\
    Integración temporal             & \multicolumn{2}{c}{Generalized-$\alpha$ ($\rho_\infty=0.9$)} \\
    \bottomrule
  \end{tabular}
\end{table}

\subsection{Parámetros del acoplamiento FE\textsuperscript{2}}

\begin{table}[htbp]
  \centering
  \caption{Configuración del acoplamiento multiescala.}
  \label{tab:ms16-fe2-params}
  \begin{tabular}{lcc}
    \toprule
    Parámetro & Valor \\
    \midrule
    Esquema de acoplamiento     & IteratedTwoWayFE2 (doble) \\
    Máx.\ iteraciones staggered & 6 \\
    Tolerancia (Frobenius)      & 3\,\% \\
    Factor de relajación        & 0.7 \\
    Paso de inicio de acoplamiento & 10 \\
    Columnas críticas           & 5 (top-N por índice de daño) \\
    Malla del submodelo         & $4\times4\times8$ Hex27 \\
    Modelo constitutivo 3D      & Ko-Bathe Concrete 3D \\
    \bottomrule
  \end{tabular}
\end{table}


% ------------------------------------------------------------------
\section{Implementación}
\label{sec:ms16-impl}

\subsection{Reasignación de grupos por rango}

Dado que la función \texttt{make\_building\_domain} genera todas las
columnas bajo un único grupo \texttt{"Columns"}, se realiza un
postproceso inmediato sobre el dominio para reasignar los grupos
según la coordenada~$z$ del nodo inferior de cada columna:

\begin{lstlisting}[language=C++,
  caption={Reasignación de grupos de columnas por rango de piso.}]
std::size_t idx = 0;
for (auto& geom : domain.elements()) {
    if (geom.physical_group() == "Columns") {
        const double z_bot = geom.point_p(0).coord(2);
        const int r = story_range(z_bot);
        geom.set_physical_group(COL_GROUPS[r]);
        elem_to_range[idx] = r;
    }
    ++idx;
}
\end{lstlisting}

Esto permite registrar tres materiales distintos en el
\texttt{StructuralModelBuilder}:

\begin{lstlisting}[language=C++,
  caption={Asignación de material por grupo de piso.}]
for (int r = 0; r < NUM_RANGES; ++r)
    builder.set_frame_material(COL_GROUPS[r],
                               col_mats[r]);
builder.set_frame_material("Beams", bm_mat);
\end{lstlisting}

\subsection{Coordinadores multiescala por rango}

Los elementos críticos se agrupan por rango de piso, y se construye
un \texttt{MultiscaleCoordinator} independiente para cada rango,
utilizando las dimensiones de sección correspondientes:

\begin{lstlisting}[language=C++,
  caption={Coordinadores por rango para submodelos con secciones distintas.}]
std::map<int, MultiscaleCoordinator> coordinators;
for (auto& [range, eids] : crit_by_range) {
    auto& coord = coordinators[range];
    for (auto eid : eids)
        coord.add_critical_element(
           extract_beam_kinematics(eid));
    coord.build_sub_models(SubModelSpec{
        .section_width  = COL_B[range],
        .section_height = COL_H[range],
        .nx=4, .ny=4, .nz=8,
        .hex_order = HexOrder::Quadratic});
}
\end{lstlisting}

\subsection{Puente C++--Python para postproceso}

Se implementó una clase utilitaria \texttt{PythonPlotter} en
\texttt{src/utils/PythonPlotter.hh} que invoca scripts de Python
directamente desde el ejecutable C++ al finalizar la simulación:

\begin{lstlisting}[language=C++,
  caption={Invocación de Python desde C++ para generación de figuras.}]
fall_n::PythonPlotter plotter(
    BASE + "scripts/falln_postprocess.py");
plotter.plot(rec_dir,
    BASE + "doc/figures/lshaped_multiscale_16/");
\end{lstlisting}

La librería de postproceso \texttt{scripts/falln\_postprocess.py}
contiene funciones reutilizables:
\texttt{plot\_roof\_displacement},
\texttt{plot\_hysteresis},
\texttt{plot\_crack\_evolution},
\texttt{plot\_interstory\_drift}, y una interfaz CLI con
\texttt{argparse}.


% ------------------------------------------------------------------
\section{Resultados esperados}
\label{sec:ms16-results}

\subsection{Desplazamientos de cubierta}

Se registran los desplazamientos de todos los nodos activos del
nivel 16 en las tres direcciones.  La historia temporal permite
identificar:
\begin{itemize}
  \item El instante de máxima respuesta ($t^* \approx 60$--$90$\,s).
  \item La amplificación debida al acoplamiento horizontal-vertical.
  \item El desfase del movimiento del suelo en las direcciones NS y EW.
\end{itemize}

La trayectoria orbital $(u_x,\,u_y)$ del nodo de cubierta más cercano
a la esquina re-entrante muestra la naturaleza bidireccional de la
demanda sobre la estructura.

\subsection{Curvas de histéresis a ambas escalas}

\begin{description}
  \item[Escala global (fibras)]
    Las fibras de acero de las columnas críticas exhiben ciclos
    completos de fluencia cíclica con endurecimiento de Menegotto-Pinto.
    Las fibras de concreto muestran degradación progresiva de rigidez
    con el modelo de Kent-Park confinado.
  \item[Escala local (submodelos 3D)]
    Los puntos de Gauss del submodelo siguen la envolvente de
    Ko-Bathe con formación de fisuras direccionales.  Se espera
    degradación más pronunciada en los submodelos del rango inferior
    (columnas de 0.50\,m, más cargadas axialmente).
\end{description}

\subsection{Evolución de fisuras en los submodelos}

Cada submodelo exporta:
\begin{itemize}
  \item Número de puntos de Gauss fisurados vs.\ tiempo.
  \item Apertura máxima de fisura vs.\ tiempo.
  \item Índice de daño máximo en la malla local.
  \item Glyphs VTK de planos de fisura para visualización en ParaView.
\end{itemize}

La comparación entre submodelos de distintos rangos permite cuantificar
el efecto de la irregularidad en altura sobre la distribución de daño.


% ------------------------------------------------------------------
\section{Discusión y análisis de validez}
\label{sec:ms16-discussion}

\subsection{Localización del daño en transiciones de sección}

La hipótesis H1 se verifica observando qué submodelos acumulan mayor
número de fisuras.  Se espera que los elementos en los pisos 5--6
(transición inferior--medio) sean los más afectados, seguidos por
los pisos 11--12 (transición medio--superior).

\subsection{Convergencia del acoplamiento FE\textsuperscript{2}}

El número de iteraciones staggered por paso ($\leq 6$) y la reducción
de la norma de Frobenius confirman la estabilidad del esquema
two-way con relajación $\omega = 0{.}7$.  Si en algún paso se alcanza
el máximo de iteraciones sin convergencia, el algoritmo continúa con
la mejor estimación disponible, registrando un aviso.

\subsection{Efecto del amortiguamiento de Rayleigh}

Con $\xi = 5\,\%$ entre $T_1 = 1{.}6$\,s y $T_3 = 0{.}4$\,s, los
modos intermedios están ligeramente sobreamortiguados.  Se recomienda
un estudio paramétrico variando $\xi$ y verificando su impacto en la
deriva inter-piso máxima.


% ------------------------------------------------------------------
\section{Estrategia de validación}
\label{sec:ms16-validation}

Se propone una estrategia de validación en tres niveles:

\begin{enumerate}
  \item \textbf{Verificación numérica (code-to-code):}
    Comparar los desplazamientos de cubierta y derivas máximas con un
    modelo equivalente en OpenSees utilizando las mismas secciones,
    materiales y registro sísmico.  El criterio de aceptación es una
    diferencia menor al 10\,\% en el pico de desplazamiento.

  \item \textbf{Validación con resultados experimentales:}
    Cotejar las curvas de histéresis de fibra y las apariciones de
    fisura con los datos de ensayos de columnas RC cíclicas publicados
    por Saatcioglu \& Ozcebe (1989) y Sheikh \& Khoury (1993).

  \item \textbf{Análisis de sensibilidad:}
    Evaluar la robustez del modelo variando:
    \begin{itemize}
      \item Factor de escala del sismo ($\alpha \in \{0.5,\, 1.0,\, 1.5\}$).
      \item Número de submodelos ($N_{\text{crit}} \in \{3,\, 5,\, 7\}$).
      \item Malla del submodelo ($4\times 4\times 8$ vs.\
            $6\times 6\times 12$).
      \item Factor de relajación ($\omega \in \{0.5,\, 0.7,\, 0.9\}$).
    \end{itemize}
\end{enumerate}

\subsection{Indicadores clave de rendimiento (KPIs)}

\begin{table}[htbp]
  \centering
  \caption{KPIs para la validación del modelo multiescala.}
  \label{tab:ms16-kpi}
  \begin{tabular}{lll}
    \toprule
    KPI & Métrica & Criterio \\
    \midrule
    Desplazamiento de cubierta & $u_{\max}$ relativo & $<10\,\%$ vs.\ referencia \\
    Deriva inter-piso máxima & $\theta_{\max}$ & $<2{.}5\,\%$ (LS) \\
    Iteraciones FE$^2$ & promedio por paso & $\leq 4$ (ideal) \\
    Convergencia Frobenius & norma relativa final & $<3\,\%$ \\
    Fisuras acumuladas & total en submodelos & monótonamente creciente \\
    \bottomrule
  \end{tabular}
\end{table}


% ------------------------------------------------------------------
\section{Conclusiones del ejemplo}
\label{sec:ms16-conclusions}

Este ejemplo demuestra que el módulo multiescala de \texttt{fall\_n}
es capaz de:
\begin{itemize}
  \item Manejar edificios de altura media (16 pisos) con
        irregularidades de sección en altu\-ra, sin modificar la API
        de \texttt{make\_building\_domain}.
  \item Ejecutar acoplamiento FE\textsuperscript{2} doble con
        múltiples rangos de sección, cada uno con su propio
        \texttt{MultiscaleCoordinator} y propiedades de material.
  \item Generar resultados de postproceso (VTK, CSV, figuras)
        de forma integrada, incluyendo la invocación de scripts
        Python desde C++.
  \item Proporcionar una base sólida para estudios paramétricos y
        de validación que se desarrollarán en capítulos futuros.
\end{itemize}
